% Vietnamese translation for OI Public Library
% Source-PDF: USACO/USACO上海交通大学暑假集训讲义54P.pdf
\documentclass[11pt]{article}
\usepackage{oipl-vn}

\title{Bài giảng tập huấn hè USACO}
\author{Ma Rong, Đại học Giao thông Thượng Hải}
\date{}

\newcommand{\problem}[1]{\subsection{#1}}
\newcommand{\iosec}[1]{\paragraph{#1}}
\newcommand{\bigO}[1]{\ensuremath{O(#1)}}

\begin{document}
\maketitle

\origtitle{Summer Training Handout, Shanghai Jiao Tong University}

\section*{Mục Lục}

\begin{itemize}
  \item Bài 1: Vét cạn và tham lam: Fair Shuttle, Fliptile, Face The Right Way.
  \item Bài 2: Bài toán ba lô: Score Inflation, Money Systems, Cow Exhibition, Space Elevator, Milk Measuring.
  \item Bài 3: Chọn lọc quy hoạch động: Apple Catching, Cheapest Palindrome, Beef McNuggets, Rebuilding Roads.
  \item Bài 4: Cây khung nhỏ nhất: Building Roads, Cheering up the Cows, Earthquake.
  \item Bài 5: Đường đi ngắn nhất: Revamping Trails, Roadblocks, Cow Jogging.
  \item Bài 6: Duyệt theo chiều rộng: Bronze Lilypad Pond, Silver Lilypad Pond, Lilypad Pond.
  \item Bài 7: Chọn lọc đề thi USACO: Moo University - Financial Aid, Team Tryouts, Emergency Pizza Order.
  \item Bài 8: Chọn lọc đề thi USACO, tiếp theo: Balanced Lineup, Ranking the Cows, Cow Traffic, Ural 1039.
  \item Bài 9: Cây đoạn.
\end{itemize}

\section*{Từ vựng thường gặp trong đề USACO}

\begin{center}
\begin{tabular}{ll}
\texttt{pasture} & đồng cỏ \\
\texttt{barn} & chuồng bò \\
\texttt{trail} & đường mòn \\
\texttt{grazing location} & nơi chăn thả \\
\texttt{intersection} & giao lộ \\
\texttt{FJ} & viết tắt của Farmer John, nông dân John \\
\texttt{Bessie} & Bessie, một con bò cái rất thông minh \\
\texttt{hay} & cỏ khô \\
\texttt{hoof} & móng bò, số nhiều là \texttt{hooves}
\end{tabular}
\end{center}

\section{Vét Cạn Và Tham Lam}

\problem{Xe đưa đón hội chợ (Fair Shuttle, USACO 2009 Feb)}

Đi chợ phiên, đổi quà và xem biểu diễn không phải việc gì khó với nông dân John, nhưng đàn bò của ông lại rất thiếu rèn luyện. Nếu phải đi chơi cả ngày ở hội chợ, chắc chắn chúng sẽ kiệt sức. Vì vậy John muốn cho bò dùng xe đưa đón trong khu hội chợ.

Tuy nhiên John không có nhiều tiền. Chiếc xe ông thuê chỉ có thể chạy một lượt theo một đường thẳng trong hội chợ, và chỉ dừng ở \(N\) địa điểm, \(1 \le N \le 20000\), được đánh số từ \(1\) đến \(N\). Đàn bò được chia thành \(K\) nhóm, \(1 \le K \le 50000\). Nhóm thứ \(i\) có \(M_i\) con, \(1 \le M_i \le N\), muốn đi từ \(S_i\) đến \(T_i\), với \(1 \le S_i < T_i \le N\).

Sức chứa xe có hạn nên có thể không chở được toàn bộ bò muốn đi; trong trường hợp đó, một phần của một nhóm cũng được phép tách ra để đi xe. Biết sức chứa xe là \(C\), \(1 \le C \le 100\), hãy giúp John lập kế hoạch thỏa mãn được nhiều bò nhất.

\iosec{Dữ liệu vào}
\begin{itemize}
  \item Dòng đầu gồm ba số nguyên \(K,N,C\), cách nhau bởi dấu cách.
  \item \(K\) dòng tiếp theo: dòng \(i+1\) chứa \(S_i,E_i,M_i\), mô tả nhóm bò thứ \(i\).
\end{itemize}

\iosec{Dữ liệu ra}
\begin{itemize}
  \item Một dòng: số bò lớn nhất có thể đi xe.
\end{itemize}

\iosec{Ví dụ}
\begin{verbatim}
shuttle.in          shuttle.out
8 15 3             10
1 5 2
13 14 1
5 8 3
8 14 2
14 15 1
9 12 1
12 15 2
4 6 1
\end{verbatim}

Xe có thể chở \(2\) con từ \(1\) đến \(5\), \(3\) con từ \(5\) đến \(8\), \(2\) con từ \(8\) đến \(14\), \(1\) con từ \(9\) đến \(12\), \(1\) con từ \(13\) đến \(14\), và \(1\) con từ \(14\) đến \(15\).

\iosec{Ý tưởng}
Tham lam: ở mọi thời điểm, luôn ưu tiên chở những con bò xuống xe sớm nhất. Bản gốc không trình bày phần chứng minh.

\problem{Lật ô (Fliptile, USACO 2007 Nov)}

John biết rằng bò thông minh hơn có thể cho nhiều sữa hơn. Ông thiết kế cho bò một trò chơi trí tuệ tên là Fliptile. Bàn chơi gồm \(M \times N\) ô, \(1 \le M,N \le 15\). Mỗi ô có hai mặt màu: một đen, một trắng.

Khi lật một ô, màu của ô đó đảo ngược. Nếu biến được tất cả các ô thành trắng thì bò thắng. Tuy nhiên móng bò rất to: khi chúng muốn lật một ô, các ô kề cạnh với ô đó cũng bị lật. Việc lật đi lật lại rất nhàm chán, nên bò muốn số lần lật là ít nhất.

Hãy xác định số lần lật nhỏ nhất và cách lật tương ứng. Nếu có nhiều nghiệm tối ưu, in nghiệm nhỏ nhất theo thứ tự từ điển. Nếu không thể hoàn thành, chỉ cần in một dòng \texttt{IMPOSSOBLE}.

\iosec{Dữ liệu vào}
\begin{itemize}
  \item Dòng đầu: hai số nguyên \(M,N\).
  \item \(M\) dòng tiếp theo: dòng \(i+1\) mô tả hàng thứ \(i\), gồm \(N\) số nguyên, \(1\) là đen và \(0\) là trắng.
\end{itemize}

\iosec{Dữ liệu ra}
\begin{itemize}
  \item \(M\) dòng, mỗi dòng \(N\) số nguyên, biểu thị số lần lật tại ô đó.
\end{itemize}

\iosec{Ví dụ}
\begin{verbatim}
fliptile.in        fliptile.out
4 4                0 0 0 0
1 0 0 1            1 0 0 1
0 1 1 0            1 0 0 1
0 1 1 0            0 0 0 0
1 0 0 1
\end{verbatim}

Còn một nghiệm khác là
\begin{verbatim}
0 1 1 0
0 0 0 0
0 0 0 0
0 1 1 0
\end{verbatim}
nhưng nghiệm này không nhỏ hơn ví dụ theo thứ tự từ điển.

\iosec{Ý tưởng}
Nếu cách lật hàng đầu tiên đã cố định thì cách lật các hàng còn lại cũng cố định. Vì vậy chỉ cần vét cạn mọi cách lật hàng đầu tiên.

\problem{Quay bò về đúng hướng (Face The Right Way, USACO 2007 Mar)}

John xếp \(N\) con bò thành một hàng, \(1 \le N \le 5000\). Phần lớn bò đều tốt và quay mặt về phía trước, nhưng một số ít lại quay về phía sau. Để đạt sự hoàn hảo, John muốn tất cả bò quay mặt về phía trước.

May mắn là John vừa mua một máy lật bò tự động. Vì là máy giảm giá, trước khi dùng ông phải đặt số bò bị lật mỗi lần là \(K\), \(1 \le K \le N\), và sau khi đặt thì không thể thay đổi. Máy chỉ có thể lật đồng thời \(K\) con bò đứng liền nhau, không được ít hơn \(K\), kể cả ở hai đầu hàng. Một con đang quay trước sẽ thành quay sau, và ngược lại.

Vì John chỉ được chọn một \(K\) không đổi, hãy tìm \(K\) sao cho số lần dùng máy ít nhất, đồng thời tính số lần ít nhất \(M\). Nếu có nhiều \(K\) thỏa mãn, in \(K\) nhỏ nhất.

\iosec{Dữ liệu vào}
\begin{itemize}
  \item Dòng đầu: số nguyên \(N\).
  \item \(N\) dòng tiếp theo: mỗi dòng là một ký tự \texttt{B} hoặc \texttt{F}, cho biết con bò thứ \(i\) quay sau hay quay trước.
\end{itemize}

\iosec{Dữ liệu ra}
\begin{itemize}
  \item Một dòng gồm hai số nguyên \(K,M\).
\end{itemize}

\iosec{Ví dụ}
\begin{verbatim}
cowturn.in         cowturn.out
7                  3 3
B
B
F
B
F
B
B
\end{verbatim}

Chọn \(K=3\), lật ba lần các đoạn \((1,2,3)\), \((3,4,5)\), \((5,6,7)\).

\iosec{Ý tưởng}
Liệt kê \(K\). Với mỗi \(K\), quét từ đầu hàng và lật tại những vị trí bắt buộc. Trạng thái của mỗi vị trí chỉ phụ thuộc vào parity số lần lật trong \(K-1\) vị trí trước nó, nên tính \(M\) trong thời gian tuyến tính. Tổng độ phức tạp là \(\bigO{n^2}\).

\section{Bài Toán Ba Lô}

Tài liệu tham khảo trong bản gốc: \url{http://xxjs.ayyz.cn/pack/Index.html}.

\subsection{Lý thuyết cơ bản}

Cho \(n\) vật có giá trị \(v_i\) và trọng lượng \(w_i\). Khi tổng trọng lượng không vượt quá, hoặc đúng bằng, \(C\), cần tìm phương án ba lô tối ưu.

Gọi \(f[i,j]\) là phương án tối ưu khi chỉ dùng \(i\) vật đầu và tổng trọng lượng không vượt quá, hoặc đúng bằng, \(j\).

Nếu yêu cầu ``không vượt quá'', điều kiện biên là
\[
  f[0,0]=f[0,1]=\cdots=f[0,C]=0.
\]
Nếu yêu cầu ``đúng bằng'', điều kiện biên là
\[
  f[0,0]=0,\qquad f[0,1]=\cdots=f[0,C]=-\infty.
\]
Công thức chuyển:
\[
  f[i,j]=\max\left(f[i-1,j], f[i-1,j-w_i]+v_i\right).
\]
Mục tiêu là \(f[n,C]\).

Các biến thể thường gặp gồm vật phẩm 0/1, vật phẩm vô hạn, vật phẩm có số lượng giới hạn, vật phẩm nhiều chi phí, vật phẩm theo nhóm, vật phẩm phụ thuộc vào vật chính, in phương án tối ưu, và in nghiệm tốt thứ \(k\).

Hai chương trình sau đều đúng, lần lượt cho ba lô 0/1 và ba lô vô hạn:
\begin{verbatim}
fillchar(f, sizeof(f), 0);          fillchar(f, sizeof(f), 0);

for i := 1 to n do                  for i := 1 to n do
  for j := m downto v[i] do           for j := v[i] to m do
    f[j] :=                              f[j] :=
      max(f[j], f[j-v[i]] + w[i]);         max(f[j], f[j-v[i]] + w[i]);

writeln(f[m]);                      writeln(f[m]);
\end{verbatim}

\problem{Lạm phát điểm số (Score Inflation, USACO 3.1)}

Trong cuộc thi USACO, thí sinh đạt điểm càng cao thì chúng ta càng vui. Ta muốn giúp thí sinh đạt nhiều điểm hơn trong cuộc thi.

Bài thi có thể phân loại theo dạng bài. Mỗi dạng có vô hạn bài, và mọi bài cùng dạng có thời gian giải cùng điểm số như nhau. Nhiệm vụ là cho biết nên chọn bao nhiêu bài thuộc mỗi dạng để thí sinh đạt tổng điểm lớn nhất trong thời gian cho phép.

Dữ liệu gồm tổng thời gian \(M\), \(1 \le M \le 10000\), và số dạng bài \(N\), \(1 \le N \le 10000\). Mỗi dòng tiếp theo gồm điểm và thời gian của một dạng bài, đều nằm trong \([1,10000]\). Mỗi dạng có thể không chọn, chọn một bài, hoặc chọn nhiều bài.

\iosec{Dữ liệu vào}
\begin{itemize}
  \item Dòng đầu: \(M,N\), tổng thời gian và số dạng bài.
  \item \(N\) dòng tiếp theo: điểm và thời gian của từng dạng.
\end{itemize}

\iosec{Dữ liệu ra}
\begin{itemize}
  \item Một dòng: điểm lớn nhất có thể đạt.
\end{itemize}

\iosec{Ví dụ}
\begin{verbatim}
inflate.in         inflate.out
300 4              605
100 60
250 120
120 100
35 20
\end{verbatim}

Chọn \(4\) bài loại hai và \(3\) bài loại bốn.

\iosec{Ý tưởng}
Đây là ba lô vô hạn điển hình. Với \(mt_i\) là thời gian và \(pt_i\) là điểm:
\[
  f[i,j]=\max\left(f[i-1,j], f[i,j-mt_i]+pt_i\right),
\]
biên \(f[0,0]=\cdots=f[0,M]=0\), mục tiêu \(f[N,M]\).

\problem{Hệ thống tiền tệ (Money Systems, USACO 2.3)}

Bò không chỉ lập chính phủ riêng mà còn xây dựng hệ thống tiền tệ riêng. Các mệnh giá của chúng rất kỳ lạ; theo thói quen, một hệ tiền thường có các mệnh giá như \(1,5,10,20\) hoặc \(25,50,100\), đôi khi thêm \(2\).

Bò muốn biết có bao nhiêu cách biểu diễn một số tiền bằng một bộ mệnh giá. Ví dụ với hệ \(\{1,2,5,10,\ldots\}\), số \(18\) có thể biểu diễn bằng \(18\) đồng \(1\), \(9\) đồng \(2\), \(8\) đồng \(2\) cộng \(2\) đồng \(1\), \(3\) đồng \(5\) cộng \(2\) cộng \(1\), v.v.

Cho các mệnh giá, hãy tính số cách biểu diễn một số tiền. Đáp án bảo đảm không vượt quá kiểu số nguyên 64 bit có dấu.

\iosec{Dữ liệu vào}
\begin{itemize}
  \item Dòng đầu: \(V,N\), với \(1 \le V \le 25\), \(1 \le N \le 10000\).
  \item Các dòng sau chứa đủ \(V\) mệnh giá, số lượng số trên mỗi dòng không cố định.
\end{itemize}

\iosec{Dữ liệu ra}
\begin{itemize}
  \item Một dòng: số cách dùng \(V\) loại tiền để biểu diễn \(N\).
\end{itemize}

\iosec{Ví dụ}
\begin{verbatim}
money.in           money.out
3 10               10
1 2 5
\end{verbatim}

\iosec{Ý tưởng}
Gọi \(c_1,c_2,\ldots,c_V\) là các mệnh giá. Gọi \(f[i,j]\) là số cách dùng \(i\) mệnh giá đầu để biểu diễn đúng \(j\):
\[
  f[0,0]=1,\quad f[0,1]=\cdots=f[0,N]=0,\quad
  f[i,j]=f[i-1,j]+f[i,j-c_i].
\]
Mục tiêu là \(f[V,N]\).

\problem{Triển lãm bò (Cow Exhibition, USACO 2003 Nov)}

Bò muốn chứng minh với xã hội rằng chúng vừa thông minh vừa thú vị. Bessie tổ chức một hội chợ triển lãm bò và đã phỏng vấn \(N\) con bò, \(1 \le N \le 100\). Mỗi con có chỉ số thông minh \(S_i\) và chỉ số thú vị \(F_i\), đều trong \([-1000,1000]\).

Bessie phải quyết định chọn những con nào tham gia. Gọi \(TS=\sum S_i\), \(TF=\sum F_i\) trên các con được chọn. Cô muốn tối đa hóa \(TS+TF\), đồng thời cả hai giá trị đều không âm. Hãy tìm tổng lớn nhất có thể.

\iosec{Dữ liệu vào}
\begin{itemize}
  \item Dòng đầu: \(N\).
  \item \(N\) dòng tiếp theo: \(S_i,F_i\).
\end{itemize}

\iosec{Dữ liệu ra}
\begin{itemize}
  \item Một dòng: giá trị lớn nhất của \(TS+TF\) với \(TS,TF\ge 0\). Nếu không có nghiệm, in \(0\).
\end{itemize}

\iosec{Ví dụ}
\begin{verbatim}
smrtfun.in         smrtfun.out
5                  8
-5 7
8 -6
6 -3
2 1
-8 -5
\end{verbatim}

Chọn bò \(1,3,4\), khi đó \(TS=3\), \(TF=5\), tổng bằng \(8\). Nếu thêm bò \(2\) thì tổng tăng lên \(10\), nhưng \(TF\) âm nên không hợp lệ.

\iosec{Ý tưởng}
Ba lô 0/1: xem \(S_i\) là dung lượng, \(S_i+F_i\) là giá trị. Cần dịch chỉ số để xử lý tổng \(S_i\) âm.

\problem{Thang máy vũ trụ (Space Elevator, USACO 2005 Mar)}

Bò chuẩn bị bay lên không gian. Để vào quỹ đạo, chúng định xây một thang máy vũ trụ khổng lồ bằng nhiều loại vật liệu. Có \(K\) loại vật liệu, \(1 \le K \le 400\). Vật liệu loại \(i\) có chiều cao \(h_i\), \(1 \le h_i \le 100\), số lượng \(c_i\), \(1 \le c_i \le 10\). Do tia vũ trụ có thể gây nguy hiểm, mỗi khối loại \(i\) khi xếp không được có độ cao vượt quá \(a_i\), \(1 \le a_i \le 40000\).

Hãy dùng các vật liệu này để dựng thang máy cao nhất có thể.

\iosec{Dữ liệu vào}
\begin{itemize}
  \item Dòng đầu: \(K\).
  \item \(K\) dòng tiếp theo: \(h_i,a_i,c_i\).
\end{itemize}

\iosec{Dữ liệu ra}
\begin{itemize}
  \item Một dòng: chiều cao lớn nhất \(H\).
\end{itemize}

\iosec{Ví dụ}
\begin{verbatim}
elevator.in        elevator.out
3                  48
7 40 3
5 23 8
2 52 6
\end{verbatim}

Từ dưới lên có thể dùng \(3\) khối loại \(2\), \(2\) khối loại \(1\), rồi \(6\) khối loại \(3\). Nếu dùng \(4\) khối loại \(2\) và \(3\) khối loại \(1\) thì không hợp lệ vì khối loại \(1\) vượt độ cao \(40\).

\iosec{Ý tưởng}
Một biến thể ba lô. Trước hết sắp các loại theo \(a_i\), rồi tách mỗi khối thành vật phẩm 0/1, sao cho \(a_1\le a_2\le\cdots\le a_n\) và \(c_i=1\) sau khi tách.
\[
  f[0,0]=\text{True},\quad f[0,j]=\text{False}\ (j>0),
\]
\[
  f[i,j]=
  \begin{cases}
    f[i-1,j]\lor f[i-1,j-h_i], & j\le a_i,\\
    \text{False}, & j>a_i.
  \end{cases}
\]
Mục tiêu là \(\max\{j\mid f[n,j]=\text{True}\}\).

\problem{Đong sữa (Milk Measuring, USACO 5.3)}

John cần đong đúng \(Q\) quart sữa thượng hạng, \(1 \le Q \le 20000\), rồi đóng vào chai gửi cho khách. Lượng sữa trong chai phải đúng như khách đặt.

John rất tiết kiệm. Để đong \(Q\) quart từ bể sữa khổng lồ, ông phải mua một số xô đo ở cửa hàng. Vì mọi xô có cùng giá, hãy tính bộ xô ít nhất cần mua. Nếu có nhiều bộ cùng ít xô nhất, John thích bộ ``nhỏ hơn'': sắp hai bộ tăng dần rồi so từ phần tử đầu; bộ có dung tích nhỏ hơn tại vị trí khác nhau đầu tiên được chọn.

Khi đong sữa, John đổ đầy xô từ bể rồi đổ vào chai. Ông không bao giờ đổ sữa từ chai ra ngoài hoặc từ xô này sang xô khác. Dữ liệu bảo đảm có ít nhất một nghiệm.

\iosec{Dữ liệu vào}
\begin{itemize}
  \item Dòng đầu: \(Q\).
  \item Dòng hai: \(P\), \(1 \le P \le 100\), số xô trong cửa hàng.
  \item \(P\) dòng tiếp theo: dung tích từng xô, từ \(1\) đến \(10000\).
\end{itemize}

\iosec{Dữ liệu ra}
\begin{itemize}
  \item Một dòng: số xô ít nhất, rồi danh sách dung tích các xô cần mua theo thứ tự tăng.
\end{itemize}

\iosec{Ví dụ}
\begin{verbatim}
milk4.in           milk4.out
16                 2 3 5
3
3
5
7
\end{verbatim}

\iosec{Ý tưởng}
Đây là ba lô vô hạn nhưng không chỉ cần giá trị tối ưu mà còn cần phương án tối ưu theo thứ tự từ điển, nên phần cài đặt khó hơn đáng kể.

\section{Chọn Lọc Quy Hoạch Động}

Dynamic Programming, viết tắt DP, không dùng chữ ``programming'' theo nghĩa lập trình máy tính, mà là ``use of a tabular solution method'', tức dùng bảng để giải bài toán.

\paragraph{Loại bài nào là bài toán quy hoạch động?}
Quy hoạch động là một phương pháp giải, không phải tên một lớp bài toán cụ thể. Khác với hình học tính toán hay lý thuyết đồ thị, quy hoạch động không có đối tượng nghiên cứu cố định. Khi nói một bài là ``bài quy hoạch động'', thường chỉ có nghĩa là bài đó có thể giải bằng quy hoạch động. Thông thường ta dùng quy hoạch động cho bài toán tối ưu; một số sách cũng gọi vài bài toán đếm cổ điển là quy hoạch động, nhưng cách gọi này không thật cần thiết.

\paragraph{Quy hoạch động và chia để trị khác nhau, liên hệ thế nào?}
Cả hai đều bắt nguồn từ quy nạp. Khác biệt là chia để trị tách bài toán ban đầu thành các bài toán con độc lập, còn quy hoạch động thường xử lý các bài toán con không độc lập và có chồng lấn.

\paragraph{Các bước giải bằng quy hoạch động}
\begin{enumerate}
  \item Đặc trưng cấu trúc nghiệm tối ưu, tức tìm bài toán con.
  \item Định nghĩa đệ quy giá trị nghiệm tối ưu, tức viết công thức chuyển.
  \item Tính giá trị nghiệm tối ưu theo chiều từ dưới lên.
  \item Dựng nghiệm tối ưu từ thông tin đã tính, bước này thường không được trình bày.
\end{enumerate}

\paragraph{Quy hoạch động có đồng nghĩa với truy hồi không?}
Không. Quy hoạch động là phương pháp giải, còn truy hồi là cách cài đặt đối lập với đệ quy. Mọi chương trình có thể cài bằng truy hồi đều có thể viết lại bằng đệ quy cộng ghi nhớ, và ngược lại. Ví dụ dãy Fibonacci \(f_n=f_{n-1}+f_{n-2}\) có thể tính từ trên xuống bằng đệ quy hoặc từ dưới lên bằng điều kiện đầu \(f_0=f_1=1\). Khi cài đặt bằng đệ quy trong bài DP, ta thường gọi là tìm kiếm có ghi nhớ.

\paragraph{Các yếu tố của quy hoạch động}
\begin{itemize}
  \item Phương trình chuyển trạng thái.
  \item Hàm mục tiêu hoặc trạng thái mục tiêu.
  \item Điều kiện biên hoặc điều kiện ban đầu.
\end{itemize}

\paragraph{Điều kiện áp dụng}
Thứ nhất là nguyên lý tối ưu: chiến lược con của một chiến lược tối ưu cũng là tối ưu. Nói cách khác, quyết định tối ưu của bài toán chỉ phụ thuộc vào quyết định tối ưu của các bài toán con; quyết định không tối ưu của bài toán con không ảnh hưởng đến lời giải. Ví dụ, nếu đường ngắn nhất từ \(A\) đến \(B\) đi qua \(C\), thì nó phải chứa đường ngắn nhất từ \(A\) đến \(C\).

Thứ hai là tính không hậu hiệu: tương lai không phụ thuộc vào quá khứ, chỉ phụ thuộc vào hiện tại. Khi trạng thái của một giai đoạn đã xác định, diễn biến về sau không bị ảnh hưởng bởi các trạng thái và quyết định trước đó.

\problem{Bắt táo (Apple Catching, USACO 2004 Nov)}

Bò thích ăn táo. John có hai cây táo đánh số \(1\) và \(2\). Bessie không với tới táo trên cây nên phải chờ chúng rơi xuống, và phải bắt trước khi táo chạm đất. Mỗi phút, một trong hai cây rơi một quả táo. Nếu Bessie đứng dưới cây đó thì bắt được.

Bessie có thể di chuyển rất nhanh giữa hai cây, nhưng mỗi phút chỉ đứng dưới một cây. Vì bò không thích vận động quá nhiều, cô chỉ sẵn lòng di chuyển tối đa \(W\) lần trong \(T\) phút, với \(1 \le T \le 1000\), \(1 \le W \le 30\). Ban đầu Bessie đứng dưới cây \(1\). Hãy tìm số táo nhiều nhất có thể bắt.

\iosec{Dữ liệu vào}
\begin{itemize}
  \item Dòng đầu: \(T,W\).
  \item \(T\) dòng tiếp theo: cây có táo rơi trong từng phút.
\end{itemize}

\iosec{Dữ liệu ra}
\begin{itemize}
  \item Một dòng: số táo lớn nhất bắt được khi di chuyển không quá \(W\) lần.
\end{itemize}

\iosec{Ví dụ}
\begin{verbatim}
bcatch.in          bcatch.out
7 2                6
2
1
1
2
2
1
1
\end{verbatim}

Bessie đứng ở cây \(1\) bắt hai quả, sang cây \(2\) bắt hai quả, rồi quay về cây \(1\) bắt hai quả cuối.

\iosec{Ý tưởng}
Quy hoạch động với chiều thứ nhất là thời gian, chiều thứ hai là số lần đã di chuyển. Từ số lần di chuyển có thể suy ra cây hiện tại.

\problem{Palindrome rẻ nhất (Cheapest Palindrome, USACO 2007 Open)}

Theo dõi từng con bò là việc khó, nên John lắp hệ thống nhận dạng tự động. Mỗi con có nhãn điện tử; khi bò đi qua máy quét, hệ thống đọc được chuỗi định danh dài \(M\), \(1 \le M \le 2000\), gồm các chữ cái thường.

Bò nghịch ngợm, đôi khi đi lùi qua máy quét. Khi đó nhãn \texttt{abcb} có thể đọc thành \texttt{abcb} hoặc \texttt{bcba}; còn \texttt{abcba} thì không có vấn đề này. John muốn biến mọi định danh thành chuỗi đọc xuôi ngược như nhau. Có thể chèn hoặc xóa ký tự ở vị trí bất kỳ; xâu rỗng cũng là palindrome.

Mỗi thao tác chèn hoặc xóa một chữ cái có chi phí riêng, từ \(0\) đến \(10000\). Cho nhãn ban đầu và chi phí chèn/xóa từng chữ cái, hãy tìm chi phí nhỏ nhất để biến nhãn thành palindrome.

\iosec{Dữ liệu vào}
\begin{itemize}
  \item Dòng đầu: \(N,M\).
  \item Dòng hai: chuỗi ban đầu dài đúng \(M\).
  \item \(N\) dòng tiếp theo: một chữ cái và hai số nguyên, lần lượt là chi phí chèn và xóa chữ cái đó.
\end{itemize}

\iosec{Dữ liệu ra}
\begin{itemize}
  \item Một dòng: chi phí nhỏ nhất.
\end{itemize}

\iosec{Ví dụ}
\begin{verbatim}
cheappal.in        cheappal.out
3 4                900
abcb
a 1000 1100
b 350 700
c 200 800
\end{verbatim}

Chèn \texttt{a} cuối chuỗi tốn \(1000\), xóa \texttt{a} đầu chuỗi tốn \(1100\), còn chèn \texttt{bcb} ở đầu tốn \(350+200+350=900\), là tối ưu.

\iosec{Ý tưởng}
Gọi \(f[i,j]\) là chi phí nhỏ nhất để biến đoạn con \(S[i..j]\) thành palindrome. Biên \(f[i,i]=0\). Chuyển:
\[
f[i,j]=\min\begin{cases}
f[i+1,j]+\operatorname{add}(S_i),\\
f[i+1,j]+\operatorname{del}(S_i),\\
f[i,j-1]+\operatorname{add}(S_j),\\
f[i,j-1]+\operatorname{del}(S_j),\\
f[i+1,j-1]\quad\text{nếu }S_i=S_j.
\end{cases}
\]
Mục tiêu là \(f[1,M]\).

\problem{McNuggets bò (Beef McNuggets, USACO 4.1)}

Đàn bò của nông dân Brown phản đối vì nghe tin McDonald's đang cân nhắc sản phẩm mới: McNuggets bò. Một chiến lược của bò là ``đóng gói kém''. Nếu chỉ có hộp \(3\), \(6\), \(10\) miếng thì không thể phục vụ khách muốn mua \(1,2,4,5,7,8,11,14,17\) miếng.

Cho số loại hộp \(N\), \(1 \le N \le 10\), và dung lượng các hộp \(i\), \(1 \le i \le 256\). In số miếng lớn nhất không thể đóng gói bằng số lượng hộp không hạn chế. Nếu mọi số đều đóng gói được hoặc các số không đóng gói được là vô hạn, in \(0\). Nếu tồn tại, số không đóng gói được lớn nhất không vượt quá \(2\cdot 10^9\).

\iosec{Dữ liệu vào}
\begin{itemize}
  \item Dòng đầu: số loại hộp \(N\).
  \item \(N\) dòng tiếp theo: số miếng mà từng loại hộp chứa.
\end{itemize}

\iosec{Dữ liệu ra}
\begin{itemize}
  \item Một dòng: số lớn nhất không thể đóng gói, hoặc \(0\).
\end{itemize}

\iosec{Ví dụ}
\begin{verbatim}
nuggets.in         nuggets.out
3                  17
3
6
10
\end{verbatim}

\iosec{Ý tưởng}
Bản thân lời giải rất đơn giản nhưng ẩn chứa một số kiến thức số học. Nếu \(\gcd\) của mọi dung lượng lớn hơn \(1\) thì có vô hạn số không biểu diễn được. Nếu \(\gcd=1\), có thể đánh dấu khả đạt theo kiểu ba lô vô hạn cho tới khi gặp một đoạn đủ dài các số liên tiếp đều đạt được.

\problem{Xây lại đường (Rebuilding Roads, USACO Feb 2002)}

Sau một trận động đất lớn, bò xây lại trang trại của John. Có \(N\) chuồng, \(1 \le N \le 150\), đánh số \(1\) đến \(N\). Vì không đủ thời gian xây nhiều đường, giữa hai chuồng bất kỳ có đúng một đường đi, tức mạng đường là một cây.

John muốn biết nếu có động đất nữa thì thiệt hại tối thiểu thế nào. Cụ thể, cần phá ít nhất bao nhiêu đường để tách được một cây con gồm đúng \(P\) chuồng, \(1 \le P \le N\), khỏi phần còn lại.

\iosec{Dữ liệu vào}
\begin{itemize}
  \item Dòng đầu: \(N,P\).
  \item \(N-1\) dòng tiếp theo: \(I,J\), cho biết \(I\) là cha của \(J\).
\end{itemize}

\iosec{Dữ liệu ra}
\begin{itemize}
  \item Một dòng: số đường ít nhất cần phá.
\end{itemize}

\iosec{Ví dụ}
\begin{verbatim}
roads.in           roads.out
11 6               2
1 2
1 3
1 4
1 5
2 6
2 7
2 8
4 9
4 10
4 11
\end{verbatim}

Cây con bị tách là \(1,2,3,6,7,8\), cần phá các cạnh \((1,4)\) và \((1,5)\).

\iosec{Ý tưởng}
DP trên cây. Gọi \(f[u,k]\) là số cạnh ít nhất cần xóa trong cây gốc \(u\) để tách một cây con kích thước \(k\) chứa \(u\). Gọi \(f^{(i)}[u,k]\) là giá trị tương tự khi chỉ xét \(i\) nhánh con đầu của \(u\). Nếu con thứ \(i\) là \(u_i\):
\[
  f[u,k]=f^{(|S(u)|)}[u,k],
\]
\[
  f^{(0)}[u,1]=0,\qquad
  f^{(i)}[u,k]=\min_{1\le j\le k-1}
  \left(f^{(i-1)}[u,j]+f[u_i,k-j]\right),
\]
đồng thời có phương án không lấy nhánh con, khi đó phải cắt thêm một cạnh. Mục tiêu:
\[
  \min_u
  \begin{cases}
    f[u,P], & u=\text{root},\\
    f[u,P]+1, & u\ne\text{root}.
  \end{cases}
\]
Bản gốc gợi ý biểu diễn cây bằng cấu trúc ``con phải, anh em trái''.

\section{Cây Khung Nhỏ Nhất}

\problem{Xây đường (Building Roads, USACO 2007 Dec)}

John vừa có vài trang trại mới và muốn nối chúng lại. Có \(N\) trang trại, \(1 \le N \le 1000\), đánh số \(1\) đến \(N\). Trang trại \(i\) có tọa độ \((X_i,Y_i)\), với \(1 \le X_i,Y_i \le 1000000\). Đã có \(M\) con đường, \(1 \le M \le 1000\). Hãy tính tổng chiều dài đường mới ít nhất cần xây để mọi trang trại liên thông.

\iosec{Dữ liệu vào}
\begin{itemize}
  \item Dòng đầu: \(N,M\).
  \item \(N\) dòng tiếp theo: \(X_i,Y_i\).
  \item \(M\) dòng tiếp theo: \(i,j\), một đường đã tồn tại giữa \(i\) và \(j\).
\end{itemize}

\iosec{Dữ liệu ra}
\begin{itemize}
  \item Một dòng: tổng chiều dài nhỏ nhất cần xây, giữ hai chữ số thập phân.
\end{itemize}

\iosec{Ví dụ}
\begin{verbatim}
roads.in           roads.out
4 1                4.00
1 1
3 1
2 3
4 3
1 4
\end{verbatim}

Xây đường giữa \(1\) và \(2\), giữa \(3\) và \(4\), mỗi đường dài \(2.00\).

\iosec{Ý tưởng}
Bài cây khung nhỏ nhất điển hình. Vì đồ thị đầy đủ theo khoảng cách Euclid là đồ thị dày, dùng Prim là phù hợp. Thực ra còn cách \(\bigO{n\log n+m}\): các cạnh có khả năng nằm trong MST thuộc tam giác Delaunay, chỉ có \(\bigO{n}\) cạnh; có thể dựng tam giác Delaunay trong \(\bigO{n\log n}\), rồi chạy Kruskal nhanh cho đồ thị thưa.

\problem{An ủi bò (Cheering up the Cows, USACO 2008 Nov)}

John trở nên rất lười và muốn bỏ càng nhiều đường càng tốt trong mạng \(P\) đường nối \(N\) đồng cỏ, \(5 \le N \le 10000\), \(N-1 \le P \le 100000\), nhưng vẫn giữ liên thông. Đồng cỏ \(i\) có một con bò. Đường thứ \(j\) nối \(S_j,E_j\) và đi hết \(L_j\) thời gian.

Bò buồn vì hệ thống giao thông bị cắt giảm. Bạn phải đến nơi ở của từng con để an ủi. Mỗi lần đến đồng cỏ \(i\), kể cả đã đến nhiều lần, bạn phải mất \(C_i\), \(1 \le C_i \le 1000\), thời gian trò chuyện. Bạn được chọn một đồng cỏ xuất phát; sau khi an ủi xong tất cả bò, phải quay lại đó ngủ qua đêm. Khi xuất phát buổi sáng và khi quay về buổi tối, bạn cũng phải trò chuyện với bò ở đồng cỏ xuất phát.

Nếu John nghe lời bạn để chọn các đường giữ lại, hãy tính thời gian ít nhất để an ủi tất cả bò.

\iosec{Dữ liệu vào}
\begin{itemize}
  \item Dòng đầu: \(N,P\).
  \item \(N\) dòng tiếp theo: \(C_i\).
  \item \(P\) dòng tiếp theo: \(S_j,E_j,L_j\).
\end{itemize}

\iosec{Dữ liệu ra}
\begin{itemize}
  \item Một dòng: thời gian ít nhất.
\end{itemize}

\iosec{Ví dụ}
\begin{verbatim}
cheer.in           cheer.out
5 7                176
10
10
20
6
30
1 2 5
2 3 5
2 4 12
3 4 17
2 5 15
3 5 6
4 5 12
\end{verbatim}

Giữ một cây phù hợp và xuất phát ở đồng cỏ \(4\), có thể đi theo thứ tự \(4,5,4,2,3,2,1,2,4\), tổng thời gian \(176\).

\iosec{Ý tưởng}
Với một cây đã chọn và gốc \(R\), lộ trình tối ưu là một vòng Euler trên cây: đi xuống từng con, xử lý xong thì quay lại cha. Một đỉnh có bậc \(D\) được thăm \(D\) lần; riêng gốc được thăm thêm một lần. Mỗi cạnh được đi đúng hai lần.

Tổng thời gian cho cây có cạnh dài \(L_i\), đỉnh bậc \(D_j\), chi phí trò chuyện \(C_j\) là
\[
  2\sum L_i+\sum D_jC_j+C_R.
\]
Chọn gốc có \(C_R\) nhỏ nhất là độc lập. Vì mỗi cạnh nối \(S,E\) đóng góp \(C_S+C_E\) vào \(\sum D_jC_j\), ta gán trọng số mới cho cạnh:
\[
  w(S,E)=2L+C_S+C_E.
\]
Khi đó cần tìm cây khung nhỏ nhất theo trọng số mới, rồi cộng \(\min C_R\). Dùng Kruskal hoặc Prim trong \(\bigO{P\log P}\).

\problem{Động đất (Earthquake, USACO 2001 Open)}

Một trận động đất phá hủy trang trại của John. John đã xây lại \(N\) trang trại, \(1 \le N \le 400\), và muốn xây đường để nối chúng. Có \(M\) đường hai chiều có thể xây, \(1 \le M \le 10000\). John muốn chi phí càng thấp càng tốt.

Đúng lúc đó, bò lập đội công trình chuyên sửa đường sau động đất. Tuy nhiên chúng chỉ làm nếu có lãi. Chúng quan tâm đến tốc độ kiếm tiền, tức tổng lợi nhuận chia cho tổng thời gian thi công. John và bò thỏa thuận tổng phí sửa chữa là \(F\), \(1 \le F \le 2000000000\). Mỗi đường có thời gian \(t\) và chi phí xây \(c\), đều trong \([1,2000000000]\). Có thể có nhiều đường giữa cùng một cặp trang trại. Bảo đảm tồn tại đường nối liên thông, nhưng chi phí có thể vượt \(F\).

Hãy chọn các đường sửa để lợi nhuận trên một đơn vị thời gian là lớn nhất.

\iosec{Dữ liệu vào}
\begin{itemize}
  \item Dòng đầu: \(N,M,F\).
  \item \(M\) dòng tiếp theo: \(i,j,c,t\), đường nối \(i,j\), thời gian \(t\), chi phí \(c\).
\end{itemize}

\iosec{Dữ liệu ra}
\begin{itemize}
  \item Một dòng: lợi nhuận đơn vị lớn nhất với bốn chữ số thập phân; nếu không có lãi, in \texttt{0.0000}.
\end{itemize}

\iosec{Ví dụ}
\begin{verbatim}
quake.in           quake.out
5 5 100            1.0625
1 2 20 5
1 3 20 5
1 4 20 5
1 5 20 5
2 3 23 1
\end{verbatim}

Chọn bốn đường cuối, tổng thời gian \(16\), tổng chi phí \(83\), lợi nhuận đơn vị là \(17/16=1.0625\).

\iosec{Ý tưởng}
Tìm kiếm tham số. Nhị phân câu hỏi: có nghiệm với lợi nhuận đơn vị vượt \(\alpha\) không? Điều này tương đương tồn tại cây khung \(T\) sao cho
\[
  \text{lợi nhuận}(T)-\alpha\cdot \text{thời gian}(T)\ge 0,
\]
tức
\[
  F-\sum_{e\in T}c_e-\alpha\sum_{e\in T}t(e)\ge 0
  \quad\Longleftrightarrow\quad
  \sum_{e\in T}\left(c_e+\alpha t(e)\right)\le F.
\]
Đặt \(w_e=c_e+\alpha t(e)\); bài kiểm tra là tìm MST theo trọng số \(w\).

\section{Đường Đi Ngắn Nhất}

Các thuật toán đường đi ngắn nhất quen thuộc về cơ bản đều là thuật toán gán nhãn. Có thể chia thành hai nhóm lớn:
\begin{itemize}
  \item Label Setting, như Dijkstra: mỗi bước xác định chắc chắn khoảng cách ngắn nhất tới một đỉnh.
  \item Label Correcting, như Bellman-Ford: khoảng cách tới mỗi đỉnh được cải thiện dần.
\end{itemize}

Dijkstra cộng heap là lựa chọn ưu tiên. Khi cài đặt, có thể dùng mẹo không cần cập nhật khóa trong heap; xem chương trình của bài Revamping Trails.

Bản gốc cũng lưu ý về SPFA: thuật toán này được công bố năm 1994, nhưng hiệu quả không cao như một số lời truyền miệng. Phân tích độ phức tạp trung bình \(\bigO{k|E|}\) với \(k\approx 2\) không có cơ sở lý thuyết rõ ràng; trường hợp xấu nhất của SPFA là \(\bigO{|V||E|}\), lớn hơn nhiều so với Dijkstra. Các liên kết tham khảo trong bản gốc:
\begin{itemize}
  \item \url{http://dantvt.spaces.live.com/blog/cns!D87988A6CAC0A480!775.entry}
  \item \url{http://dantvt.spaces.live.com/blog/cns!D87988A6CAC0A480!790.entry}
  \item \url{http://hi.baidu.com/winterlegend/blog/item/8d9b16d801c861e038012f5e.html}
\end{itemize}
Dữ liệu bài Revamping Trails có quy mô khoảng \(100000\), khá thích hợp để so sánh.

\problem{Nâng cấp đường mòn (Revamping Trails, USACO 2009 Feb)}

Mỗi ngày John phải kiểm tra bò trong chuồng. Ông cần đi từ chuồng \(1\) tới chuồng \(N\) qua đường gần nhất. Có \(N\) chuồng và \(M\) đường hai chiều, \(1 \le M \le 50000\), bảo đảm chuồng \(1\) nối được tới chuồng \(N\). Đường \(i\) nối \(P1_i,P2_i\) và mất \(T_i\), \(1 \le T_i \le 1000000\), thời gian.

John muốn nâng cấp một số đường để giảm thời gian đi lại. Ông chỉ có thể nâng cấp tối đa \(K\) đường, \(1 \le K \le 20\), và đường đã nâng cấp có thời gian đi bằng \(0\). Hãy chọn cách nâng cấp để đường từ \(1\) tới \(N\) ngắn nhất.

\iosec{Dữ liệu vào}
\begin{itemize}
  \item Dòng đầu: \(N,M,K\).
  \item \(M\) dòng tiếp theo: \(P1_i,P2_i,T_i\).
\end{itemize}

\iosec{Dữ liệu ra}
\begin{itemize}
  \item Một dòng: độ dài đường ngắn nhất sau khi nâng cấp không quá \(K\) đường.
\end{itemize}

\iosec{Ví dụ}
\begin{verbatim}
revamp.in          revamp.out
4 4 1              1
1 2 10
2 4 10
1 3 1
3 4 100
\end{verbatim}

Nâng cấp đường \(3\to4\), đường \(1\to3\to4\) có tổng thời gian \(1\).

\iosec{Ý tưởng}
Độ dài \(K\)-sửa của một đường đi là tổng độ dài sau khi bỏ qua \(K\) cạnh dài nhất của đường đi đó. Cần tìm đường từ \(1\) đến \(N\) có độ dài \(K\)-sửa nhỏ nhất; nó không nhất thiết là đường đi ngắn nhất ban đầu rồi sửa.

Từ đồ thị \((V,E)\), dựng đồ thị mới với đỉnh \(V\times\{0,1,\ldots,K\}\). Với mỗi cạnh \(e=(v_1,v_2)\) dài \(L\), thêm cạnh \((v_1,k)\to(v_2,k)\) dài \(L\) cho mọi \(k\), và cạnh \((v_1,k)\to(v_2,k+1)\) dài \(0\) cho \(0\le k<K\); làm tương tự chiều ngược vì đường hai chiều. Thêm cạnh \((v,k)\to(v,k+1)\) dài \(0\). Khi đó đường ngắn nhất từ \((1,0)\) đến \((N,K)\) chính là đáp án. Đồ thị mới có \(\bigO{VK}\) đỉnh và \(\bigO{EK}\) cạnh, chạy Dijkstra trong \(\bigO{EK\log V}\).

\problem{Rào cản đường (Roadblocks, USACO 2006 Nov)}

Bessie chuyển đến một trang trại nhỏ và thỉnh thoảng về thăm bạn. Vì phong cảnh trên đường đẹp, cô không muốn đi quá nhanh, nên chọn đường ngắn thứ hai thay vì đường ngắn nhất. Giả sử đường ngắn thứ hai luôn tồn tại.

Có \(R\) đường hai chiều, \(1 \le R \le 100000\), và \(N\) nút, \(1 \le N \le 5000\). Bessie đi từ \(1\) tới \(N\). Đường ngắn thứ hai có thể dùng lại cạnh của đường ngắn nhất và được phép có chu trình, tức có thể đi qua một đỉnh hoặc cạnh nhiều lần. Đường ngắn thứ hai phải dài hơn nghiêm ngặt đường ngắn nhất; nếu có hai đường khác nhau cùng ngắn nhất, chúng đều không tính là thứ hai.

\iosec{Dữ liệu vào}
\begin{itemize}
  \item Dòng đầu: \(N,R\).
  \item \(R\) dòng tiếp theo: \(A,B,D\), đường nối \(A,B\) dài \(D\).
\end{itemize}

\iosec{Dữ liệu ra}
\begin{itemize}
  \item Một dòng: độ dài đường ngắn thứ hai từ \(1\) đến \(N\).
\end{itemize}

\iosec{Ví dụ}
\begin{verbatim}
block.in           block.out
4 4                450
1 2 100
2 4 200
2 3 250
3 4 100
\end{verbatim}

Đường ngắn nhất là \(1\to2\to4\), dài \(300\). Đường thứ hai là \(1\to2\to3\to4\), dài \(450\).

\iosec{Ý tưởng}
Gọi \(d(p,k)\) là độ dài đường ngắn thứ \(k\) tới đỉnh \(p\). Ta chỉ cần xét \(d(p,1)\) và \(d(p,2)\). Nếu một đường dùng trạng thái \(d(i,3)\), thay phần tới \(i\) bằng hai đường ứng với \(d(i,1)\), \(d(i,2)\) sẽ tạo các đường rẻ hơn, mâu thuẫn với việc đường ban đầu thuộc hai đường tốt nhất.

Vì mọi cạnh có độ dài dương, trạng thái có khoảng cách nhỏ nhất trong hàng đợi không thể bị giảm nữa. Dùng hàng đợi ưu tiên; mỗi lần lấy trạng thái nhỏ nhất, nếu đỉnh của nó chưa được mở rộng hai lần thì mở rộng sang mọi láng giềng. Không cần mở rộng cùng một đỉnh quá hai lần. Có nhiều nhất \(2V\) lần mở rộng và hàng đợi chứa cỡ \(2E\) trạng thái, thời gian \(\bigO{E\log E}\).

\problem{Bò chạy chậm (Cow Jogging, USACO 2008 Mar)}

Bessie phải giảm cân nên mỗi tuần chạy bộ vài lần giữa chuồng và ao. Cô chỉ muốn chạy đường xuống dốc từ chuồng đến ao, rồi đi bộ về. Các nút đánh số \(1\) đến \(N\), \(1 \le N \le 1000\). Nếu \(X>Y\), nút \(X\) cao hơn \(Y\), nên đường xuống dốc đi từ \(X\) đến \(Y\). Chuồng của Bessie là nút \(N\), điểm cao nhất; ao là nút \(1\), điểm thấp nhất.

Sau một tuần, Bessie chán tuyến đường đơn điệu và muốn có \(K\) lựa chọn khác nhau, \(1 \le K \le 100\). Để không quá mệt, cô muốn đó là \(K\) đường ngắn nhất từ chuồng đến ao. Mỗi đường có dạng \((X_i,Y_i,D_i)\), \(1 \le Y_i<X_i\le N\), dài \(D_i\).

\iosec{Dữ liệu vào}
\begin{itemize}
  \item Dòng đầu: \(N,M,K\).
  \item \(M\) dòng tiếp theo: \(X_i,Y_i,D_i\).
\end{itemize}

\iosec{Dữ liệu ra}
\begin{itemize}
  \item \(K\) dòng: dòng \(i\) là độ dài đường ngắn thứ \(i\), hoặc \(-1\) nếu không tồn tại. Nếu nhiều đường có cùng độ dài, in lặp lại độ dài đó.
\end{itemize}

\iosec{Ví dụ}
\begin{verbatim}
jogging.in         jogging.out
5 8 7              1
5 4 1              2
5 3 1              2
5 2 1              3
5 1 1              6
4 3 4              7
3 1 1              -1
3 2 1
2 1 1
\end{verbatim}

Sáu đường theo thứ tự độ dài gồm \(5\to1\), \(5\to3\to1\), \(5\to2\to1\), \(5\to3\to2\to1\), \(5\to4\to3\to2\to1\), và một đường dài hơn.

\iosec{Ý tưởng}
Đồ thị có hướng không chu trình và có thứ tự tự nhiên theo chỉ số đỉnh, nên dùng DP. Giả sử đã biết \(K\) đường ngắn nhất bắt đầu từ mọi đỉnh lớn hơn \(i\) và đi xuống ao. Để tính cho \(i\), xét mọi cạnh ra khỏi \(i\); với mỗi cạnh, ghép cạnh đó vào trước từng đường tốt nhất của đỉnh đích. Chọn \(K\) giá trị nhỏ nhất trong các danh sách đã sắp. Nếu đỉnh có bậc ra \(d\), bước trộn tốn \(dK\); tổng thời gian \(\bigO{MK}\).

\section{Duyệt Theo Chiều Rộng}

\problem{Ao sen đồng (Bronze Lilypad Pond, USACO 2007 Feb)}

Để bò giải trí và rèn luyện, John xây một ao hình chữ nhật chia thành \(M\) hàng \(N\) cột, \(1 \le M,N \le 30\). Một số ô là lá sen chắc chắn, một số là đá, còn lại là nước. Bessie tập múa ballet trên một lá sen và muốn nhảy tới một lá sen khác. Cô chỉ được nhảy từ lá sen sang lá sen, không được xuống nước hay đá.

Bước nhảy giống quân mã: mỗi lần đi ngang \(M_1\) ô rồi dọc \(M_2\) ô, hoặc dọc \(M_1\) rồi ngang \(M_2\), với \(1 \le M_1,M_2 \le 30\), \(M_1\ne M_2\). Tối đa có tám hướng. Cho bố cục ao và độ dài bước nhảy, hãy tính số bước ít nhất từ điểm đầu đến điểm cuối. Dữ liệu bảo đảm đến được đích.

\iosec{Dữ liệu vào}
\begin{itemize}
  \item Dòng đầu: \(M,N,M_1,M_2\).
  \item \(M\) dòng tiếp theo: mỗi dòng \(N\) số, trong đó \(0\) là nước, \(1\) là lá sen, \(2\) là đá, \(3\) là điểm đầu, \(4\) là điểm cuối.
\end{itemize}

\iosec{Dữ liệu ra}
\begin{itemize}
  \item Một dòng: số bước ít nhất.
\end{itemize}

\iosec{Ví dụ}
\begin{verbatim}
bronlily.in        bronlily.out
4 5 1 2            2
1 0 1 0 1
3 0 2 0 4
0 1 2 0 0
0 0 0 1 0
\end{verbatim}

Bessie nhảy tới lá sen ở hàng một cột ba, rồi nhảy đến đích.

\iosec{Ý tưởng}
Biến thể đơn giản của bài quân mã, dùng BFS hoặc Dijkstra.

\problem{Ao sen bạc (Silver Lilypad Pond, USACO 2007 Feb)}

Bài toán giống ao sen đồng, nhưng bước nhảy cố định như quân mã thường: một ô theo một chiều và hai ô theo chiều kia. John nhận thấy đôi khi Bessie không tới được đích vì thiếu lá sen, nên muốn thêm một số lá sen, không được đặt lên đá. John muốn thêm ít lá sen nhất; trong số đó, cần số bước ít nhất; sau cùng, cần số đường đi có số lá sen thêm ít nhất và số bước ít nhất.

\iosec{Dữ liệu vào}
\begin{itemize}
  \item Dòng đầu: \(M,N\).
  \item \(M\) dòng tiếp theo: ma trận trạng thái, \(0\) nước, \(1\) lá sen, \(2\) đá, \(3\) điểm đầu, \(4\) điểm cuối.
\end{itemize}

\iosec{Dữ liệu ra}
\begin{itemize}
  \item Dòng một: số lá sen ít nhất cần thêm, hoặc \(-1\) nếu vô nghiệm.
  \item Dòng hai: với số lá sen thêm ít nhất, số bước ít nhất; không in nếu dòng một là \(-1\).
  \item Dòng ba: số đường đi đạt hai tiêu chí trên; không in nếu dòng một là \(-1\).
\end{itemize}

\iosec{Ví dụ}
\begin{verbatim}
silvlily.in        silvlily.out
4 8                2
00010000           6
00000201           2
00000400
30000010
\end{verbatim}

Cần thêm ít nhất hai lá sen, Bessie cần ít nhất sáu bước và có hai đường đi như vậy.

\iosec{Ý tưởng}
Xem đây là bài đường đi ngắn nhất. Ở mỗi ô lưu hai giá trị: cần thêm bao nhiêu lá sen để tới đó và đã nhảy bao nhiêu bước. Trước hết tối thiểu hóa số lá sen, sau đó tối thiểu hóa số bước. Đi vào một lá sen sẵn có tốn một bước; đi vào ô trống tốn một lá sen cộng một bước.

Dùng BFS để cập nhật. Nếu một bước nhảy đưa tới trạng thái có cùng số lá sen và số bước như đã biết, cộng số đường đi. Nếu tốt hơn, thay số đường đi và đưa ô vào hàng đợi. Nếu kém hơn thì bỏ qua.

\problem{Ao sen vàng (Lilypad Pond, USACO 2007 Feb)}

Bài toán giống Silver Lilypad Pond, nhưng yêu cầu tìm số lá sen ít nhất cần thêm và số cách đặt các lá sen đó.

\iosec{Dữ liệu vào}
\begin{itemize}
  \item Dòng đầu: \(M,N\).
  \item \(M\) dòng tiếp theo: ma trận trạng thái.
\end{itemize}

\iosec{Dữ liệu ra}
\begin{itemize}
  \item Dòng một: số lá sen ít nhất cần thêm, hoặc \(-1\) nếu vô nghiệm.
  \item Dòng hai: số cách đặt, bảo đảm không vượt quá số nguyên 64 bit có dấu; không in nếu dòng một là \(-1\).
\end{itemize}

\iosec{Ví dụ}
\begin{verbatim}
lilypad.in         lilypad.out
4 5                2
10000              3
30000
00200
00040
\end{verbatim}

Cần ít nhất hai lá sen và có ba cách đặt.

\iosec{Ý tưởng}
Xem lưới là đồ thị, các bước nhảy là cạnh. Bài toán trở thành: trong đồ thị có một số đỉnh đã đánh dấu, cần đánh dấu thêm ít đỉnh nhất để tồn tại đường đi từ \(s\) tới \(t\).

Trước hết co mỗi thành phần liên thông gồm các đỉnh đã đánh dấu thành một đỉnh, vì đi trong thành phần đó không tốn chi phí. Sau đó gán trọng số để biến bài toán thành tìm đường ngắn nhất và đếm số đường ngắn nhất. Một cạnh giữa hai đỉnh chưa đánh dấu có chi phí \(1\), vì cần đánh dấu thêm một đỉnh. Khi đi qua đỉnh đã đánh dấu, cần ``nhảy qua'' nó: nếu \(a,b\) chưa đánh dấu, \(v_1,v_2\) đã đánh dấu, và tồn tại các cạnh \(a\to v_1\), \(a\to v_2\), \(v_1\to b\), \(v_2\to b\), thì phải có cạnh trực tiếp chi phí \(1\) từ \(a\) tới \(b\) để tránh đếm cùng tập đỉnh chưa đánh dấu nhiều lần.

Số lá sen ít nhất cần thêm bằng khoảng cách ngắn nhất từ \(s\) tới \(t\) trừ \(1\), vì cạnh cuối vào \(t\) là miễn phí. Để đếm, BFS từ \(s\), rồi xét các đỉnh theo thứ tự khoảng cách và cộng số đường đi từ các láng giềng ở khoảng cách \(d-1\).

\section{Chọn Lọc Đề Thi USACO}

\problem{Đại học Moo Moo: học bổng (Moo University - Financial Aid, USACO 2004 Mar)}

Bessie thấy con người có nhiều đại học nhưng bò thì không, nên cô và bạn bè lập ra Đại học Moo Moo. Để tuyển sinh, họ phát minh bài kiểm tra năng lực học thuật bò, viết tắt CSAT; điểm của mỗi bò là một số nguyên từ \(1\) đến \(2000000000\).

Học phí ở Đại học Moo Moo rất đắt, nhiều bò cần xin học bổng từ \(1\) đến \(10000\). Không có học bổng chính phủ cho bò, nên ngân sách phải lấy từ quỹ hữu hạn của trường, tổng quỹ là \(F\), \(0 \le F \le 2000000000\).

Trường chỉ có \(N\) phòng học, \(1 \le N \le 19999\), \(N\) lẻ, trong khi có \(C\) con bò nộp đơn, \(N \le C \le 100000\). Bessie muốn nhận \(N\) con và làm cho trung vị điểm CSAT của chúng cao nhất. Trung vị là số ở giữa sau khi sắp xếp.

\iosec{Dữ liệu vào}
\begin{itemize}
  \item Dòng đầu: \(N,C,F\).
  \item \(C\) dòng tiếp theo: điểm CSAT và số học bổng mong muốn của từng bò.
\end{itemize}

\iosec{Dữ liệu ra}
\begin{itemize}
  \item Một dòng: trung vị lớn nhất có thể; nếu quỹ không đủ hỗ trợ \(N\) con bò, in \(-1\).
\end{itemize}

\iosec{Ví dụ}
\begin{verbatim}
finance.in         finance.out
3 5 70             35
30 25
50 21
20 20
5 18
35 30
\end{verbatim}

Nhận các bò có điểm \(5,35,50\), trung vị \(35\), tổng học bổng \(18+30+21=69\le70\).

\iosec{Ý tưởng}
Liệt kê ứng viên làm trung vị. Hai phía của trung vị chọn \((N-1)/2\) con rẻ nhất. Việc tìm chi phí rẻ nhất cho phía trái và phải có thể làm động bằng heap: heap chứa các chi phí đã chọn, với phần tử đắt nhất ở gốc. Khi thêm một con, nếu rẻ hơn gốc thì thay gốc. Độ phức tạp \(\bigO{C\log N}\).

\problem{Đại học Moo Moo: tuyển đội (Moo University - Team Tryouts, USACO 2004 Mar)}

Có \(N\) con bò đăng ký tuyển vào đội thể dục của Đại học Moo Moo. Mỗi con có chiều cao và chiều rộng là số nguyên dương không quá \(10000\). Cần chọn đội nhiều bò nhất sao cho với mỗi bò trong đội,
\[
  A(H-h)+B(W-w)\le C,
\]
trong đó \(H,W\) là chiều cao và chiều rộng của bò đó, \(h,w\) là chiều cao và chiều rộng nhỏ nhất trong đội, còn \(A,B,C\) là các hằng số nguyên dương không quá \(10000\).

\iosec{Dữ liệu vào}
\begin{itemize}
  \item Dòng đầu: \(N\).
  \item Dòng hai: \(A,B,C\).
  \item \(N\) dòng tiếp theo: chiều cao và chiều rộng của từng bò.
\end{itemize}

\iosec{Dữ liệu ra}
\begin{itemize}
  \item Một dòng: số bò lớn nhất trong đội.
\end{itemize}

\iosec{Ví dụ}
\begin{verbatim}
tryout.in          tryout.out
8                  5
1 2 4
5 1
3 2
2 3
2 1
7 2
6 4
5 1
4 3
\end{verbatim}

Các bò \(1,2,3,4,7\) có thể tạo thành một đội hợp lệ và không có đội lớn hơn.

\iosec{Ý tưởng}
Cách đơn giản là liệt kê mọi chiều cao nhỏ nhất và chiều rộng nhỏ nhất, rồi đếm mọi bò thỏa điều kiện; cách này \(\bigO{n^3}\), quá chậm với \(N=1000\).

Cải tiến bằng cách cố định \(minh\), rồi duyệt các khả năng \(minw\) tăng dần. Khi \(minw\) tăng, tập điểm thỏa ràng buộc tuyến tính chỉ lớn thêm. Nếu sắp các điểm theo \(Ah+Bw\), các điểm sẽ trở nên hợp lệ theo đúng thứ tự đó. Thuật toán:
\begin{itemize}
  \item Duyệt \(minh\).
  \item Duyệt \(minw\) tăng dần.
  \item Thêm các điểm mới hợp lệ vào tập và cập nhật đáp án.
  \item Loại điểm tương ứng với \(minw\) khỏi tập.
\end{itemize}
Độ phức tạp \(\bigO{n^2}\), đủ nhanh. Bản gốc dẫn thêm một \href{http://blog.doglog.net/2009/03/2008-moo-university-team-tryouts/}{bài viết trên doglog}.

\problem{Đại học Moo Moo: đặt pizza khẩn cấp (Moo University - Emergency Pizza Order, USACO 2004 Mar)}

Căng tin Đại học Moo Moo hết cỏ khô, phải lập tức đặt pizza cho \(C\) con bò, \(1 \le C \le 1000\). Cửa hàng Pizza Farm có pizza cỡ lớn vừa đủ cho một con bò. Họ sẵn sàng làm pizza cho từng con nhưng đơn hàng quá lớn nên đặt ba hạn chế:
\begin{itemize}
  \item Có \(T\) loại rau phủ, \(1 \le T \le 30\); mỗi pizza phải dùng đúng \(K\) loại, \(1 \le K \le T\).
  \item Một pizza không được dùng lặp cùng một loại phủ.
  \item Hai pizza bất kỳ không được có tập phủ hoàn toàn giống nhau.
\end{itemize}
Các loại phủ đánh số từ \(1\) đến \(T\). Mỗi con bò có một số loại phủ nó thích; pizza của nó phải có ít nhất một loại phủ nó thích. Hãy xác định nhiều nhất bao nhiêu bò có thể được phục vụ.

\iosec{Dữ liệu vào}
\begin{itemize}
  \item Dòng đầu: \(C,T,K\).
  \item \(C\) dòng tiếp theo: số loại phủ bò đó thích, rồi danh sách các loại phủ.
\end{itemize}

\iosec{Dữ liệu ra}
\begin{itemize}
  \item Một dòng: số bò tối đa được ăn pizza.
\end{itemize}

\iosec{Ví dụ}
\begin{verbatim}
pizza.in           pizza.out
3 2 1              2
2 2 1
1 1
1 2
\end{verbatim}

Chỉ có hai loại pizza khác nhau, nên không thể thỏa mãn cả ba con bò.

\iosec{Ý tưởng}
Bản chất là bài ghép cặp: ghép nhiều bò nhất với các pizza hợp lệ. Số pizza có thể rất lớn, nên sinh dần thay vì xây toàn bộ danh sách. Ghép tối ưu có thể làm bằng luồng mạng.

Cài đặt nhanh hơn bằng cách biểu diễn tập phủ bằng các bit trong số nguyên 32 bit thay vì mảng boolean tối đa \(30\) phần tử. Khi đó kiểm tra pizza có thỏa một con bò không chỉ cần vài phép bit. Trường hợp xấu gần \(\bigO{n^3}\), nhưng thực tế thường tốt hơn.

\section{Chọn Lọc Đề Thi USACO, Tiếp Theo}

\problem{Cân bằng vàng (Balanced Lineup, USACO 2007 Mar)}

Cho một dãy \(N\) con bò, mỗi con có \(K\) đặc tính được mã hóa bằng một số thập phân dài \(K\) bit. Nếu bit cuối là \(1\), con bò có đặc tính \(1\); nếu bit đầu là \(1\), nó có đặc tính \(K\). Cần tìm đoạn liên tiếp dài nhất sao cho mỗi đặc tính xuất hiện cùng số lần.

\iosec{Dữ liệu vào}
\begin{itemize}
  \item Dòng đầu: \(N,K\).
  \item \(N\) dòng tiếp theo: số thập phân biểu diễn \(K\) bit đặc tính của con bò thứ \(i\).
\end{itemize}

\iosec{Dữ liệu ra}
\begin{itemize}
  \item Một dòng: độ dài đoạn cân bằng lớn nhất.
\end{itemize}

\iosec{Ví dụ}
\begin{verbatim}
lineup.in          lineup.out
7 3                4
7
6
7
2
1
4
2
\end{verbatim}

Chọn bò \(3\) đến \(6\), độ dài \(4\), mỗi đặc tính xuất hiện đúng hai lần.

\iosec{Ý tưởng}
Xem dữ liệu là ma trận \(0/1\) kích thước \(N\times K\), gọi \(S\) là ma trận tổng tiền tố kích thước \((N+1)\times K\), với \(S[0,c]=0\), \(S[r,c]=\sum_{i\le r}A[i,c]\). Cần tìm hai hàng xa nhất \(0\le r_1<r_2\le N\) sao cho với mọi cột \(c\),
\[
  S[r_2,c]-S[r_1,c]=S[r_2,1]-S[r_1,1].
\]
Tương đương
\[
  S[r_2,c]-S[r_2,1]=S[r_1,c]-S[r_1,1].
\]
Đặt \(S'[r,c]=S[r,c]-S[r,1]\). Bài toán trở thành tìm hai vector hàng bằng nhau trong dãy \(S'\) với khoảng cách lớn nhất. Có thể sắp vector hoặc dùng bảng băm.

\problem{Xếp hạng bò (Ranking the Cows, USACO 2007 Mar)}

Mỗi con bò của John có sản lượng sữa khác nhau. John muốn xếp hạng bò theo sản lượng từ cao xuống thấp. Mỗi lần ông chỉ có thể so sánh hai con. Đã biết \(M\) kết quả so sánh, \(1 \le M \le 10000\). Để hoàn thành thứ hạng, ông muốn lập danh sách các cặp còn phải so sánh sao cho danh sách ngắn nhất; gọi độ dài là \(C\). Hãy tìm \(C\) nhỏ nhất.

\iosec{Dữ liệu vào}
\begin{itemize}
  \item Dòng đầu: \(N,M\).
  \item \(M\) dòng tiếp theo: \(X,Y\), nghĩa là bò \(X\) có sản lượng cao hơn bò \(Y\).
\end{itemize}

\iosec{Dữ liệu ra}
\begin{itemize}
  \item Một dòng: giá trị nhỏ nhất của \(C\).
\end{itemize}

\iosec{Ví dụ}
\begin{verbatim}
ranking.in         ranking.out
5 5                3
2 1
1 5
2 3
1 4
3 4
\end{verbatim}

Biết \(2>1>5\) và \(2>3>4\), nên \(2\) đứng đầu. Vẫn cần so \(1\) với \(3\), \(4\) với \(5\), và tùy kết quả còn cần một so sánh nữa, nên tối thiểu ba câu hỏi.

\iosec{Ý tưởng}
Nếu từ các kết quả đã biết không suy ra được thứ tự giữa \(X\) và \(Y\), gọi \((X,Y)\) là một cặp chưa có thứ tự. Số cặp như vậy \(C'\) chính là \(C\). Tính bao đóng bắc cầu của quan hệ lớn hơn, rồi đếm các cặp không so sánh được.

\problem{Giao thông bò (Cow Traffic, USACO 2007 Mar)}

Do số bò tăng mạnh, các đường về chuồng bị quá tải trong giờ vắt sữa. John muốn tìm con đường đông nhất để cải tạo. Có \(M\) đường một chiều, \(1 \le M \le 50000\), và \(N\) giao lộ, \(1 \le N \le 5000\), đánh số \(1\) đến \(N\). Chuồng ở giao lộ \(N\). Mỗi đường đi từ giao lộ có số nhỏ hơn tới số lớn hơn, nên đồ thị không có chu trình; mọi đường đều dẫn về chuồng. Có thể có nhiều đường giữa cùng một cặp giao lộ.

Trong giờ cao điểm, bò đi từ các điểm chăn thả tới chuồng. Các giao lộ chỉ có cạnh đi ra, không có cạnh đi vào, là điểm chăn thả. Mỗi đường đi từ một điểm chăn thả tới chuồng có một con bò khác nhau đi qua. Hãy tìm số bò đi qua con đường đông nhất. Kết quả không vượt quá số nguyên 32 bit có dấu.

\iosec{Dữ liệu vào}
\begin{itemize}
  \item Dòng đầu: \(N,M\).
  \item \(M\) dòng tiếp theo: hai đầu của một đường.
\end{itemize}

\iosec{Dữ liệu ra}
\begin{itemize}
  \item Một dòng: số đường đi qua cạnh đông nhất.
\end{itemize}

\iosec{Ví dụ}
\begin{verbatim}
traffic.in         traffic.out
7 7                4
1 3
3 4
3 5
4 6
2 3
5 6
6 7
\end{verbatim}

Cạnh \((6,7)\) nằm trên bốn đường đi từ điểm chăn thả tới chuồng.

\iosec{Ý tưởng}
Với cạnh \((u,v)\),
\[
  \#\text{đường đi qua }(u,v)
  =
  \#\text{đường từ các nguồn tới }u
  \times
  \#\text{đường từ }v\text{ tới }N.
\]
Tính hai đại lượng này bằng DP trên DAG.

\problem{Tiệc kỷ niệm trường (Ural 1039)}

Hiệu trưởng Đại học Ural chuẩn bị tiệc kỷ niệm \(80\) năm. Nhân viên trong trường có cấp bậc khác nhau, tạo thành một cây quan hệ nhân sự gốc là hiệu trưởng. Mỗi nhân viên có mã số duy nhất từ \(1\) đến \(N\), và có một giá trị vui vẻ nếu tham gia tiệc. Để mọi khách đều thoải mái, hiệu trưởng quy định một nhân viên và cấp trên trực tiếp của người đó không được cùng tham gia.

Hãy lập danh sách khách mời sao cho tổng độ vui vẻ là lớn nhất.

\iosec{Dữ liệu vào}
Dòng đầu là \(N\), \(1 \le N \le 6000\). Sau đó là \(N\) dòng, mỗi dòng là độ vui vẻ của một nhân viên, nằm trong \([-128,127]\). Tiếp theo là cây quan hệ nhân sự; mỗi dòng có \(L,K\), nghĩa là nhân viên \(K\) là cấp trên trực tiếp của nhân viên \(L\). Danh sách kết thúc bằng \texttt{0 0}.

\iosec{Dữ liệu ra}
In tổng độ vui vẻ lớn nhất của các khách mời.

\iosec{Ví dụ}
\begin{verbatim}
input              output
7                  5
1
1
1
1
1
1
1
1 3
2 3
6 4
7 4
4 5
3 5
0 0
\end{verbatim}

\iosec{Ý tưởng}
Bản gốc để trống phần lời giải.

\section{Cây Đoạn}

Bản PDF chỉ có tiêu đề bài giảng này, không có nội dung tiếp theo trong phần trích xuất.

\end{document}
